% The two simplices at (4,3): the weights, and the leverage scores through the
% substitution s_c -> 1 - s_c.  Both are 3-simplices inside R^4, drawn through
% the affine identification of the standard simplex with R^3 that carries the
% vertex e_c to the c-th vertex of the regular tetrahedron of
% figures/tetrahedron.tex,
%   x_1 = p_1 + p_2 - p_3 - p_4
%   x_2 = p_1 - p_2 + p_3 - p_4
%   x_3 = p_1 - p_2 - p_3 + p_4        (barycentre -> origin),
% followed by the orthographic projection of R^3 with azimuth 15 and elevation
% 70 degrees, scaled by 1.35cm -- the camera and the scale of
% figures/tetrahedron.tex, so the tetrahedron there and these two are the same
% drawing:
%   X = 1.35(-0.2588 x_1 + 0.9659 x_2)
%   Y = 1.35(-0.9077 x_1 - 0.2432 x_2 + 0.3420 x_3)
% The right panel repeats the left, shifted 7.7cm in X.  The leverage-score
% polytope is a simplex because the corank m - k is one; at corank two and above
% it is a hypersimplex, and this drawing does not apply.
%
% The path drawn is t = (1-3u, u, u, u) for 0 < u <= 1/4, carried at the chart
% P = I_4 - J_4/4 of the tetrahedron design: s_c = P_cc = 3/4 at every u, and
% l_c = s_c/t_c is 3/(4(1-3u)) at c = 1 and 3/(4u) at c = 2, 3, 4, all equal to
% 3 at u = 1/4.  In the identification above the path is the segment from the
% origin to the vertex e_1, since (1-3u)g_1 + u(g_2 + g_3 + g_4) = (1-4u)g_1.
\begin{tikzpicture}[
  x=1cm, y=1cm,
  vec/.style={-{Latex[length=2mm,width=1.5mm]}, line width=0.75pt},
  frame/.style={densely dotted, line width=0.3pt, black!35},
  lead/.style={line width=0.3pt, black!55},
  edge/.style={line width=0.45pt, black!65},
  hidden/.style={line width=0.45pt, black!45, dashed, dash pattern=on 1.4pt off 1.2pt},
  every node/.style={font=\small}]

% ---- one simplex, vertices e_1..e_4, barycentre at the origin -----------
% projected vertices:  e_1 -> ( 1, 1, 1) -> ( 0.9546,-1.0920)
%                      e_2 -> ( 1,-1,-1) -> (-1.6534,-1.3588)
%                      e_3 -> (-1, 1,-1) -> ( 1.6534, 0.4354)
%                      e_4 -> (-1,-1, 1) -> (-0.9546, 2.0154)
% e_2--e_3 is the far diagonal of the silhouette quadrilateral, hence hidden;
% the four dotted medians locate the barycentre against the vertices.
\newcommand{\simplexframe}{%
  \draw[hidden] (-1.6534,-1.3588)--( 1.6534, 0.4354);
  \draw[edge]   ( 0.9546,-1.0920)--(-1.6534,-1.3588)
                ( 0.9546,-1.0920)--( 1.6534, 0.4354)
                ( 0.9546,-1.0920)--(-0.9546, 2.0154)
                (-1.6534,-1.3588)--(-0.9546, 2.0154)
                ( 1.6534, 0.4354)--(-0.9546, 2.0154);
  \draw[frame]  (0,0)--( 0.9546,-1.0920) (0,0)--(-1.6534,-1.3588)
                (0,0)--( 1.6534, 0.4354) (0,0)--(-0.9546, 2.0154);
  \fill (0,0) circle (1.6pt);}

% ======================= left: the weights ==============================
\simplexframe
\node[anchor=south] at (0,2.24) {the weights};
\draw[lead] (-0.13,0)--(-1.38,0);
\node[anchor=east, align=right, font=\footnotesize] at (-1.42,0)
  {$t_c=\tfrac14$\\[1pt] {\scriptsize at $u=\tfrac14$}};
% the escape path t = (1-3u, u, u, u) is the median from the barycentre to e_1
\draw[vec] (0,0)--( 0.9546,-1.0920);
\node[anchor=north, font=\footnotesize] at (0.9546,-1.30) {$t=(1,0,0,0)$};
\node[anchor=north, align=center, font=\footnotesize] at (0,-1.90)
  {$t_c\ge0$, \ $\sum_c t_c=1$\\ $t\in\Delta^{3}$};

% ======================= middle: what each coordinate does ==============
\node[align=center, font=\footnotesize] at (3.85,0.34)
  {the chart $P$ held fixed\\[1pt]
   $t=(1-3u,u,u,u)$, \ $u\to0$\\[5pt]
   $s_c=P_{cc}$\\[5pt]
   $\ell_c=s_c/t_c$\\[1pt]
   {\scriptsize singular only at $t_c=0$}\\[5pt]
   $\ell_1\to\tfrac34$\\[1pt]
   $\ell_2=\ell_3=\ell_4=\tfrac{3}{4u}\to\infty$};

% ======================= right: the leverage scores =====================
\begin{scope}[shift={(7.7,0)}]
  \simplexframe
  \node[anchor=south] at (0,2.24) {the leverage scores};
  \draw[lead] (0.13,0)--(1.66,0);
  \node[anchor=west, align=left, font=\footnotesize] at (1.70,0)
    {$s_c=\tfrac34$\\[1pt] {\scriptsize at every $u$}};
  \node[anchor=north, font=\footnotesize] at (0.9546,-1.30) {$s=(0,1,1,1)$};
  \node[anchor=north, align=center, font=\footnotesize] at (0,-1.90)
    {$0\le s_c\le1$, \ $\sum_c s_c=3$\\ $1-s_c\in\Delta^{3}$};
\end{scope}

\end{tikzpicture}
